\chapter[ПРОЕКТИРОВАНИЕ ВЫЧИСЛИТЕЛЬНОГО ЯДРА]{ПРОЕКТИРОВАНИЕ И РЕАЛИЗАЦИЯ ВЫЧИСЛИТЕЛЬНОГО ЯДРА}

\section{Общая архитектура программного комплекса на языке \texttt{Rust}}

Программный комплекс реализован по принципу
«одно вычислительное ядро и несколько каналов доставки результата».
Все термодинамические вычисления сосредоточены в \texttt{if97\_core}.
Пользовательские приложения и сервисный контур
используют единый программный интерфейс ядра и не дублируют формулы IF97.

Такое разделение решает две инженерные задачи.
Первая задача --- совпадение результата во всех каналах доставки.
Вторая задача --- локализация математической логики в одном модуле,
что упрощает сопровождение и верификацию.

Для обмена данными применяется общий слой DTO \texttt{if97\_app\_api}.
Он фиксирует контракт входов и выходов между компонентами.
Реализация на \texttt{Rust} обеспечивает строгую типизацию,
контролируемую обработку ошибок и высокую производительность
без потери надёжности~\cite{rust_book}.

\section{Структура и модули библиотеки \texttt{if97\_core}}

Внутренняя организация \texttt{if97\_core}
построена как набор слоёв с явным разделением ответственности.
Декомпозиция соответствует практикам предметно-ориентированного проектирования (DDD):
доменная модель отделена от вычислительной логики,
а операции ядра реализованы как детерминированные процедуры без внутреннего состояния.

Типовая структура ядра включает следующие функциональные блоки:
\begin{itemize}
\item доменный модуль \texttt{domain} --- типы предметной области, единицы измерения, ошибки, итоговое состояние воды;
\item вычислительный модуль \texttt{engine} --- валидация входа, классификация области/региона и запуск плана вычисления;
\item региональный модуль \texttt{regions} --- уравнения IF97 для регионов 1--5, обратные зависимости и численные процедуры для неаналитических постановок;
\item модуль топологии \texttt{topology} --- фазовые границы, линия насыщения (регион~4), граница B23 и правила выбора региона;
\item публичный фасад \texttt{calculator} --- интерфейс \texttt{If97}, который принимает типизированный вход и возвращает унифицированный результат.
\end{itemize}

Отдельно выделена подсистема констант и таблиц коэффициентов IF97.
Она подключается к региональным моделям как неизменяемые данные.
Для задач визуализации используется интерфейс насыщения,
который строит границы купола без лишних вычислений двухфазной смеси.

\section{Использование системы типов для обеспечения размерности величин}

В ядре используется типобезопасное представление физических величин.
Давление, температура, энтальпия и другие параметры
передаются через newtype-обёртки.
Например, \texttt{MegaPascal} и \texttt{Kelvin}.

Такая типизация решает практические задачи.
Она снижает риск смешения единиц измерения.
Она делает программный интерфейс читаемым на уровне сигнатур.
Она упрощает централизованную валидацию диапазонов.

Результат расчёта представлен структурой \texttt{WaterState}.
Она содержит входные и вычисленные свойства:
\begin{itemize}
\item давление $p$ и температура $T$;
\item удельный объём $v$ и плотность $\rho$;
\item энтальпия $h$, энтропия $s$ и внутренняя энергия $u$;
\item изобарная теплоёмкость $c_p$ и скорость звука $w$;
\item номер региона IF97, в котором выполнялся расчёт.
\end{itemize}

Унифицированный формат результата
обеспечивает единое представление состояния
для всех маршрутов и всех внешних потребителей ядра.

\section{Механизм маршрутизации расчётов и классификации регионов}

Так как IF97 использует региональные уравнения,
основой вычислительного процесса является
корректное определение региона и выбор модели
~\cite{iapws_if97_2012}.
В ядре это реализовано маршрутизатором,
который анализирует вход и выбирает план расчёта.

Публичный фасад \texttt{If97} поддерживает несколько основных маршрутов:
\begin{itemize}
\item прямой маршрут $pT$ --- расчёт состояния по давлению и температуре;
\item обратные маршруты $ph$ и $ps$ --- расчёт по давлению и энтальпии или по давлению и энтропии;
\item маршрут $px$ --- расчёт в двухфазной области по давлению и степени сухости $x$;
\item маршрут $\rho T$ --- численное расширение, в котором давление подбирается методом секущих, после чего используется стандартный маршрут $pT$;
\item специализированный маршрут метастабильного пара --- режим, ограниченный областью применимости (расширение региона~2).
\end{itemize}

Внутри каждого маршрута применяется единая схема.
Сначала выполняются валидация и проверка физической допустимости входа.
Далее определяется регион с учётом топологии IF97,
линии насыщения и границы B23.
После этого выбирается вычислительная процедура:
прямая формула, аналитическая обратная зависимость
или численный решатель.

Для пар $(p, h)$ и $(p, s)$
используются аналитические обратные формулы IF97 там,
где они заданы стандартом.
Это позволяет исключить итерации на типовых участках области.
Для региона~3 дополнительно используются зависимости IAPWS,
включая релиз $v(p, T)$~\cite{iapws_vpt_region3}.

При обратных расчётах в околокритической зоне
метод Ньютона--Рафсона оказался недостаточно устойчивым
из-за деградации Якобиана.
В рабочей реализации используется двумерный решатель
Левенберга--Марквардта с сеточным поиском
и демпфированием шага.
Дополнительно реализованы защита от некорректных фазовых переходов
и выделен специализированный маршрут метастабильного пара.

\section{Оптимизация: генерация таблиц коэффициентов и организация констант}

В IF97 большинство формул представлено рядами с коэффициентами.
Для исключения ошибок ручного ввода
коэффициенты хранятся в CSV и генерируются в статический код на этапе сборки.

\texttt{build.rs} читает CSV-таблицы,
сортирует и нормализует данные
и генерирует \texttt{Rust}-модуль со статическими массивами.
Сгенерированный код подключается через \texttt{include!}
и становится частью итогового бинарного файла.

Данный подход обеспечивает следующие преимущества:
\begin{itemize}
\item отсутствие операций ввода-вывода во время выполнения при обращении к коэффициентам;
\item воспроизводимость сборки и предсказуемое поведение при развёртывании;
\item ускорение вычислений за счёт обращения к константам;
\item снижение риска ошибок из-за несовместимых форматов или окружения.
\end{itemize}

Для снижения вычислительных затрат
в горячих участках применяются
кэширование степеней $\pi$ и $\tau$
и аккуратное суммирование рядов (например, через \texttt{.fold()}).

\section{Стратегия обработки ошибок и валидация входных данных}

Для инженерного инструмента важно
не только вычислить свойства,
но и корректно обработать выход за область применимости
и некорректные комбинации входов.
Поэтому в \texttt{if97\_core}
используются централизованная валидация
и единая модель предметных ошибок.

На этапе валидации проверяются:
\begin{itemize}
\item принадлежность давления и температуры допустимым диапазонам IF97;
\item положительность плотности;
\item диапазон степени сухости $x \in [0; 1]$;
\item специальные ограничения для метастабильного пара и обратных маршрутов.
\end{itemize}

Ошибки рассматриваются как часть предметной логики
и возвращаются вызывающей стороне в диагностируемом виде
(\texttt{OutOfBounds}, \texttt{PhaseBoundaryError}, \texttt{InvalidInput}).
За счёт строгой типизации и проверок
исключаются неявные «тихие» состояния,
включая бесконтрольное распространение \texttt{NaN}.
Единая модель ошибок даёт одинаковое поведение
в \texttt{FLTK},
в оболочке на базе \texttt{Yew}, \texttt{WebAssembly} и \texttt{Tauri}
и в \texttt{gRPC}-сервисе.

\section{Верификация ядра на тестовых наборах данных IAPWS}

Верификация ядра включает модульные и интеграционные тесты,
а также сравнение с контрольными значениями.
Контрольные точки основаны на материалах IF97
и покрывают прямые и обратные маршруты~\cite{iapws_if97_2012}.

Для автоматической проверки используются:
\begin{itemize}
\item модульные тесты региональных формул, границ областей и инвариантов на реперных точках;
\item интеграционные тесты маршрутов $pT$, $ph$, $ps$, $px$ и $\rho T$ через публичный фасад \texttt{If97};
\item сценарии сверки по CSV-таблицам контрольных значений (регрессионная проверка).
\end{itemize}

Помимо корректности, проверяется производительность.
Отдельные сценарии оценивают
скорость пакетных расчётов и стоимость построения купола насыщения.
В актуальном прогоне из \texttt{TESTS.md}
подтверждены 18 модульных тестов,
1 интеграционный CSV-тест
и 1 doctest для \texttt{if97\_core}.
На уровне модульных тестов покрываются следующие аспекты:
классификация региона,
границы B23 и 2bc,
прямые и обратные уравнения регионов 1--5,
метастабильная ветвь для метастабильного пара
и непрерывность свойств на ключевых переходах.
Интеграционный CSV-тест
прогоняет 242 контрольные строки
через фасад \texttt{If97::*}
и сравнивает результаты с эталонными значениями.
Для нагрузочного сценария \texttt{load\_stress\_core\_csv}
зафиксировано $2\,009$ строк/с
(\texttt{test}-профиль, одиночный запуск).
Для многопоточного вычислительного сценария \texttt{if97\_loadtest}
при 10 прогонах (\texttt{profile=debug}, 8 потоков)
получено среднее $5\,451{,}6$ строк/с
с 95\% доверительным интервалом $[5\,053{,}88;\ 5\,849{,}32]$.
Сценарии измерения производительности и нагрузочные сценарии с атрибутом \texttt{\#[ignore]}
не смешиваются с функциональной верификацией
и запускаются отдельно как разделы замеров производительности.
Внутренний тест построения купола
по-прежнему даёт время порядка десятков миллисекунд на диаграмму.
